Thus @@M0@@ is the coefficient of @@M1@@ in the series:
\sgaStableTarget{sga6.e05.eq.5-6-1}@@M2@@

\sgaStableTarget{sga6.e05.prop.5-7}\textbf{Proposition 5.7.} \uline{Let} @@M3@@ \uline{be a binomial ring}, @@M4@@ \uline{an augmented} @@M5@@-@@M6@@-\uline{algebra}, \uline{and} @@M7@@ \uline{an element of} @@M8@@ \uline{such that} @@M9@@ \uline{for} @@M10@@. \uline{Then for every} @@M11@@ \uline{and every} @@M12@@, \uline{we have}:
\sgaStableTarget{sga6.e05.eq.5-7-1}@@M13@@

Let @@M14@@ be the @@M15@@-@@M16@@-algebra generated by generators @@M17@@, subject to the relations @@M18@@ for @@M19@@. We define the augmentation homomorphism @@M20@@ by @@M21@@. We know that @@M22@@ is the ring of polynomials with coefficients in @@M23@@ in the indeterminates @@M24@@, with @@M25@@, and @@M26@@, with @@M27@@, and that the @@M28@@-filtration on @@M29@@ is identical to the weight filtration, assigning the weight @@M30@@ to @@M31@@, @@M32@@ (4.12 and 4.15).

Put @@M33@@, @@M34@@. We shall first prove 5.7 for @@M35@@ and @@M36@@. By the preceding discussion, (5.7.1) is equivalent to:
@@M37@@
Now @@M38@@; this relation can therefore be written:
@@M39@@
by (5.5.4). This relation does hold, since @@M40@@ has zero augmentation.

By 4.7, there exists a @@M41@@-@@M42@@-homomorphism @@M43@@ compatible with the augmentations, sending @@M44@@ to @@M45@@, and @@M46@@ to @@M47@@, hence @@M48@@ to @@M49@@ and @@M50@@
